\chapter{Conception et modélisation}
\label{ch:conception}

Les diagrammes présentés dans ce chapitre ont été produits à partir du code final et non
l'inverse : les classes, les attributs et les cardinalités ont été relevés dans les entités
JPA, et les enchaînements du diagramme de séquence dans le corps des méthodes concernées.
Ils décrivent donc le système tel qu'il est, y compris lorsque cela révèle un choix
discutable.

\section{Architecture générale}

L'application adopte une architecture trois tiers classique, mais dont la couche de
sécurité présente une particularité : elle est traversée à chaque requête, et non
seulement à la connexion. La figure~\ref{fig:archi} en donne la vue d'ensemble.

\begin{figure}[H]
\centering
\includegraphics[width=\textwidth,height=0.76\textheight,keepaspectratio]{\figpath/uml/02-architecture.png}
\caption{Architecture applicative en trois tiers}
\label{fig:archi}
\end{figure}

Le client est une application monopage React communiquant exclusivement en JSON avec le
serveur. Un intercepteur Axios attache le jeton d'authentification à chaque appel et
intercepte les réponses \code{401} pour ramener l'utilisateur à l'écran de connexion.

Le serveur est organisé en quatre niveaux. Le filtre \code{JwtAuthenticationFilter}
s'exécute avant tout traitement métier ; les contrôleurs REST portent les règles
d'autorisation ; les services concentrent la logique métier et les transactions ; les
repositories Spring~Data assurent la persistance. Deux services transversaux se
distinguent : \code{AuditService}, qui écrit dans une transaction indépendante pour que la
trace subsiste même si l'opération auditée échoue, et \code{FileStorageService}, qui isole
le stockage disque derrière une abstraction unique.

L'infrastructure comprend une base PostgreSQL dont le schéma est créé et versionné par
Flyway au démarrage, un serveur SMTP pour l'envoi des identifiants et des notifications,
et un répertoire de stockage pour les fichiers déposés. L'interprétation à retenir de
cette figure est que le chemin de sécurité n'est pas contournable : aucune flèche ne relie
directement le client aux contrôleurs.

\section{Diagramme de cas d'utilisation}

Le diagramme de cas d'utilisation a été présenté au chapitre précédent
(figure~\ref{fig:usecase}) car il constitue le résultat direct de l'analyse des acteurs. Il
est rappelé ici que sa construction est \emph{descendante depuis le code} : chaque
association acteur--cas d'utilisation est adossée à une annotation d'autorisation
existante.

\section{Diagramme de classes}

Le modèle comporte 30 entités persistantes. Les représenter sur une seule planche aurait
produit une figure illisible en format A4 ; le diagramme a donc été scindé en quatre vues
cohérentes correspondant aux trois couches identifiées lors de l'analyse. Seize de ces trente
entités --- les racines métier --- héritent d'une classe abstraite \code{BaseEntity} portant
l'identifiant UUID, les métadonnées d'audit JPA (dates et auteurs de création et de
modification), le numéro de version optimiste et les deux champs de suppression logique.
Les quatorze autres --- tables d'historique, tables de liaison et journal d'audit --- portent
leur propre identifiant et ne connaissent pas la suppression logique : on n'y ajoute que
des lignes, on n'en retire jamais.

\subsection{Identité et gouvernance}

La figure~\ref{fig:classIdent} présente la vue centrale du modèle, celle où se lit la
séparation exigée par le cahier des charges.

\begin{figure}[H]
\centering
\includegraphics[width=\textwidth,height=0.78\textheight,keepaspectratio]{\figpath/uml/03-class-identite-gouvernance.png}
\caption{Diagramme de classes --- identité, membres et gouvernance}
\label{fig:classIdent}
\end{figure}

Trois observations structurent la lecture de cette figure.

D'abord, \code{User} et \code{Member} sont deux entités distinctes reliées par une
association \emph{un à zéro-ou-un}. Un \code{User} est un compte d'authentification ; un
\code{Member} est un profil scientifique. La cardinalité \emph{zéro-ou-un} n'est pas un
artefact : elle autorise l'existence d'un compte purement technique, l'administrateur
système, qui n'est pas un chercheur du laboratoire.

Ensuite, trois associations partent de \code{Member} vers des tables d'historique :
\code{MemberStatusHistory} pour le statut, \code{AffiliationHistory} pour les périodes
d'appartenance et \code{Mandate} pour les responsabilités. Une quatrième,
\code{UserRoleHistory}, part de \code{User} et enregistre le rôle applicatif. Cette
disposition matérialise la séparation : le statut, le mandat et le rôle sont historisés
\emph{séparément} parce que ce sont trois concepts distincts qui évoluent indépendamment.

Enfin, la note attachée à \code{User} porte la règle la plus importante du modèle :
l'attribut \code{role} est dérivé par \code{UserRoleSyncService} à partir du statut et des
mandats actifs, et n'est jamais saisi manuellement. Un lecteur du seul schéma relationnel
ne pourrait pas le deviner --- d'où sa présence explicite sur le diagramme.

\subsection{Activité scientifique}

La figure~\ref{fig:classActiv} regroupe les entités qui portent la production du
laboratoire.

\begin{figure}[H]
\centering
\includegraphics[width=\textwidth,keepaspectratio]{\figpath/uml/04-class-activite-scientifique.png}
\caption{Diagramme de classes --- activité scientifique}
\label{fig:classActiv}
\end{figure}

Deux choix de modélisation y sont visibles. \code{PublicationAuthor} est une classe
d'association entre \code{Publication} et \code{Member} : elle porte un attribut
\code{authorOrder} car l'ordre des signataires d'un article est une information
scientifique significative qu'une simple relation plusieurs-à-plusieurs perdrait. De même,
\code{EventOrganizer} et \code{EventParticipant} sont distinctes : organiser un événement
et s'y inscrire ne sont pas la même chose et ne relèvent pas des mêmes droits.

L'attribut \code{editionNumber} de \code{Event} répond à l'exigence d'archivage par édition
du cahier des charges : deux éditions successives d'un même colloque sont deux lignes
autonomes reliées par leur titre, et non un enregistrement écrasé.

\subsection{Ressources du laboratoire}

La figure~\ref{fig:classRess} couvre les documents, le matériel et les partenariats.

\begin{figure}[H]
\centering
\includegraphics[width=0.97\textwidth,keepaspectratio]{\figpath/uml/05-class-ressources.png}
\caption{Diagramme de classes --- documents, matériel et partenariats}
\label{fig:classRess}
\end{figure}

La relation \code{DocumentEntity}~\code{1..*}~\code{DocumentVersion} traduit l'exigence
d'historisation documentaire : un document n'est jamais remplacé, une version lui est
ajoutée. La note portée sur \code{Equipment} rappelle l'invariant de stock
--- \code{availableQuantity} diminue à la distribution et remonte au retour --- dont la
vérification fait l'objet de tests dédiés (§\ref{sec:testsequip}).

\subsection{Traçabilité et pilotage}

La figure~\ref{fig:classTrace} rassemble les entités transverses.

\begin{figure}[H]
\centering
\includegraphics[width=\textwidth,keepaspectratio]{\figpath/uml/06-class-tracabilite.png}
\caption{Diagramme de classes --- traçabilité, notifications et pilotage}
\label{fig:classTrace}
\end{figure}

\code{AuditLog} présente une particularité assumée, signalée par une note : cette entité
n'a \emph{aucune} association JPA vers les entités auditées. Elle mémorise un couple
\code{entityType}~/~\code{entityId} sous forme de chaînes. Ce couplage volontairement
faible permet de tracer n'importe quelle entité, y compris une entité archivée, sans que
la suppression logique de la cible ne rende la trace inexploitable. C'est un compromis :
on renonce à l'intégrité référentielle pour garantir la pérennité de la trace.

\section{Diagramme de séquence}

Le scénario retenu est l'attribution d'un mandat de directeur. C'est l'opération la plus
délicate du système : elle met en jeu deux membres, clôt un mandat, en ouvre un autre et
recalcule deux rôles applicatifs, le tout sans qu'aucun rôle ne soit saisi. La
figure~\ref{fig:seqMandat} en donne le déroulement exact, relevé dans
\code{MandateService.assign()} et \code{UserRoleSyncService}.

\begin{landscape}
\begin{figure}[H]
\centering
\makebox[\textwidth][c]{%
  \includegraphics[width=1.42\textwidth,keepaspectratio]{\figpath/uml/07-sequence-attribution-mandat.png}}
\caption{Diagramme de séquence --- attribution d'un mandat de directeur}
\label{fig:seqMandat}
\end{figure}
\end{landscape}

Le scénario se lit en deux temps, matérialisés par les deux blocs du diagramme.

Le premier temps assure l'unicité du titulaire. Le service recherche le mandat de directeur
actif ; s'il en existe un, il le \emph{clôt} en positionnant sa date de fin et son
indicateur d'activité, sans jamais le supprimer --- la note de droite le rappelle. Il
demande ensuite au service de synchronisation de recalculer le rôle de l'ancien titulaire.
Celui-ci, n'ayant plus de mandat actif, retrouve le rôle dérivé de son seul statut, par
exemple \code{PERMANENT\_MEMBER}. Le point important est que la perte des permissions de
direction est immédiate et automatique : elle ne dépend d'aucune action d'administration.

Le second temps ouvre le mandat du nouveau titulaire et déclenche le même recalcul. La note
centrale explicite la règle de priorité appliquée par \code{computeEffectiveRole} : un
statut gelé (retraité ou ancien membre) l'emporte sur tout, faute de quoi le mandat le plus
élevé gagne, selon l'ordre directeur $>$ vice-directeur $>$ chef d'équipe.

Chaque changement ouvre une ligne dans l'historique des rôles avec sa cause
(\code{MANDATE\_CHANGE}), et l'opération complète est journalisée. La note finale résume
l'invariant que ce scénario doit préserver : le statut des deux membres est inchangé, le
mandat est clos pour l'un et ouvert pour l'autre, et seul le rôle --- donnée dérivée --- a
été recalculé.

\section{Diagramme d'états}

Le cycle de vie d'un membre, de la candidature publique au départ ou à la retraite,
est présenté en figure~\ref{fig:etats}. Les états correspondent exactement aux valeurs des
énumérations \code{MemberStatus} et \code{RequestStatus} ; aucun état intermédiaire n'a été
ajouté pour la beauté du schéma.

\begin{figure}[H]
\centering
\includegraphics[width=\textwidth,height=0.74\textheight,keepaspectratio]{\figpath/uml/08-etats-cycle-vie-membre.png}
\caption{Diagramme d'états --- cycle de vie d'un membre}
\label{fig:etats}
\end{figure}

Trois enseignements se dégagent de cette figure.

Le composite \emph{Actif} regroupe les trois statuts entre lesquels un membre peut
librement circuler par un changement de statut. Ces transitions sont réversibles et toutes
historisées.

Les transitions vers \code{RETIRED} et \code{FORMER} sont les seules à produire un effet de
bord sur la sécurité : elles gèlent le compte et clôturent les mandats en cours. La note
attachée rappelle que ce gel est vérifié à chaque requête par le filtre d'authentification,
et non uniquement à la connexion --- un jeton encore valide ne permet donc pas de
contourner une mise à la retraite prononcée entre-temps.

Les transitions de retour, déclenchées par \code{returnToLab()}, n'annulent pas le départ :
elles ouvrent une \emph{nouvelle} période d'affiliation. C'est la traduction directe de la
règle §4.3 du cahier des charges, selon laquelle un membre peut quitter le laboratoire et
revenir sans que l'épisode précédent ne disparaisse.

\section{Modèle de données}

Le schéma relationnel compte 30 tables créées par trois migrations Flyway successives. Le
tableau~\ref{tab:migrations} en donne la répartition.

\begin{table}[H]
\centering
\caption{Migrations Flyway et contenu du schéma}
\label{tab:migrations}
\small
\begin{tabularx}{\textwidth}{@{}p{1.6cm}p{4.6cm}X@{}}
\toprule
\textbf{Version} & \textbf{Objet} & \textbf{Contenu} \\
\midrule
V1 & Schéma initial & 28 tables : identité, membres, gouvernance, activité scientifique, ressources, notifications, audit \\
V2 & Historique des rôles et identité du laboratoire & \code{user\_role\_history}, \code{laboratory} \\
V3 & Source de réconciliation & Extension de la contrainte \code{CHECK} sur \code{user\_role\_history.source} pour admettre la valeur \code{RECONCILIATION} \\
\bottomrule
\end{tabularx}
\end{table}

Le schéma répond aux exigences de normalisation du document de spécification : clés
primaires en UUID, clés étrangères explicites, index sur les colonnes de recherche
fréquente, contraintes \code{CHECK} sur les colonnes portant des énumérations, et tables
d'historique et d'audit dédiées. L'ORM est configuré en mode \code{validate} : Hibernate
vérifie la concordance entre les entités et le schéma au démarrage, mais ne le modifie
jamais. Toute évolution passe donc obligatoirement par une migration versionnée --- ce qui
rend le schéma reproductible d'un environnement à l'autre.

\section{Architecture de sécurité}

La sécurité repose sur trois principes de conception, détaillés au
chapitre~\ref{ch:securite} et introduits ici parce qu'ils relèvent de l'architecture.

\paragraph{Authentification sans état.} Aucune session serveur n'est conservée. Le client
présente un jeton signé à chaque requête. Cela simplifie le déploiement mais implique que
la révocation ne peut pas reposer sur la destruction d'une session.

\paragraph{Autorité de la base plutôt que du jeton.} C'est le choix le plus important. Le
jeton transporte bien un rôle, mais ce rôle \emph{n'est pas utilisé} pour autoriser : à
chaque requête, le filtre recharge l'utilisateur depuis la base et reconstruit ses
habilitations. Un changement de rôle ou une mise à la retraite prend donc effet à la
requête suivante, sans attendre l'expiration du jeton. Le coût est une lecture
supplémentaire par requête ; le bénéfice est qu'un jeton ne peut pas survivre à la
révocation des droits qu'il transportait.

\paragraph{Autorisation au plus près du métier.} Les règles d'accès ne sont pas centralisées
dans une table de correspondance URL~/~rôle, mais déclarées par annotation sur chaque
méthode de contrôleur. La règle est ainsi lisible à côté du code qu'elle protège, ce qui
réduit le risque de divergence lors d'une évolution.
